%-----------------------------------------------------------------------
\subsection{Instance: Mechanistic Interpretability (Circuit Discovery)}
\label{sec:instance-mech-interp}
%-----------------------------------------------------------------------

Mechanistic interpretability (MI) has emerged as a central paradigm for
AI safety research, aiming to reverse-engineer neural networks into
human-understandable circuits~\citep{elhage2022toy,conmy2023automated,anthropic2025circuit}.
We instantiate Theorem~\ref{thm:universal} for circuit discovery, where
the fundamental obstruction is that functionally equivalent networks admit
multiple valid circuit decompositions.

\paragraph{The explanation system.}
Define the 6-tuple $\mathcal{S}_{\mathrm{circuit}} =
(\Params,\, \sH,\, Y,\, \mathsf{obs},\, \mathsf{exp},\, \incomp)$ as follows.

\begin{itemize}
  \item $\Params$ is the set of neural network weight configurations
    (architecture plus trained parameters).
  \item $\sH$ is the set of \emph{circuit decompositions}: subgraphs of the
    computational graph together with an interpretation of which components
    implement which sub-computations (e.g., ``heads 1.3 and 2.1 form an
    induction circuit'').
  \item $Y$ is the set of input--output functions: each $y \in Y$ is the
    function $f_\theta : \mathcal{X} \to \mathcal{Y}$ computed by the network
    on the domain of interest.
  \item $\mathsf{obs} : \Params \to Y$ maps each weight configuration
    $\theta$ to the function it computes.
  \item $\mathsf{exp} : \Params \to \sH$ maps each network to its circuit
    decomposition---the internal computational mechanism identified by a
    circuit discovery method.
  \item $\incomp \subseteq \sH \times \sH$ is the \emph{component-attribution
    incompatibility}: $(h, h') \in \incomp$ if $h$ and $h'$ attribute the same
    computation to disjoint or contradictory subsets of network components (e.g.,
    $h$ says the computation is implemented by heads in layer 2, while $h'$ says
    it is implemented by heads in layer 4).
\end{itemize}

\paragraph{The Rashomon property: circuit multiplicity.}
\citet{meloux2025everything} provide striking empirical evidence that the
Rashomon property holds for circuit discovery.  For a simple XOR task, they
identified \textbf{85 distinct valid circuits} with \emph{zero circuit error}
(each circuit perfectly reproduces the network's behavior on the task).
Moreover, each circuit admitted an average of \textbf{535.8 valid
interpretations}---different assignments of computational roles to
components, all equally faithful to the network's input--output behavior.

Complementary evidence comes from \citet{paulo2025sae}, who show that
sparse autoencoders trained on the \emph{same} data learn \emph{different}
features depending on the random seed, demonstrating that even the
feature-level decomposition underlying circuit identification is not unique.

\begin{proposition}[Circuit Rashomon Property]
\label{prop:circuit-rashomon}
$\mathcal{S}_{\mathrm{circuit}}$ has the Rashomon property: there exist
networks $\theta_1, \theta_2 \in \Params$ with
$\mathsf{obs}(\theta_1) = \mathsf{obs}(\theta_2)$ (identical I/O behavior)
yet $(\mathsf{exp}(\theta_1), \mathsf{exp}(\theta_2)) \in \incomp$
(incompatible circuit decompositions).
\end{proposition}

\begin{proof}
Consider any pair of distinct valid circuits from the 85 found by
\citet{meloux2025everything}: by construction, both achieve zero circuit
error (same I/O behavior on the task), yet they attribute the computation
to different subsets of network components.
\end{proof}

\paragraph{Instance: mechanistic interpretability impossibility.}

\begin{theorem}[Mechanistic Interpretability Impossibility]
\label{thm:mech-interp-impossibility}
The explanation system $\mathcal{S}_{\mathrm{circuit}}$ has the Rashomon
property.  Therefore, by Theorem~\ref{thm:universal}, no circuit
explanation method $\mathsf{exp}$ can simultaneously be:
\begin{enumerate}
  \item \textbf{Faithful}: $\mathsf{exp}(\theta)$ reflects the actual
    internal circuitry of network $\theta$;
  \item \textbf{Stable}: $\mathsf{exp}(\theta_1) = \mathsf{exp}(\theta_2)$
    whenever $\mathsf{obs}(\theta_1) = \mathsf{obs}(\theta_2)$ (functionally
    equivalent networks receive the same circuit explanation); and
  \item \textbf{Decisive}: $\mathsf{exp}$ commits to a single circuit
    decomposition (rather than reporting the space of valid circuits).
\end{enumerate}
\end{theorem}

\begin{proof}
By Proposition~\ref{prop:circuit-rashomon},
$\mathcal{S}_{\mathrm{circuit}}$ has the Rashomon property.
Apply Theorem~\ref{thm:universal}.
\end{proof}

This is verified in Lean~4 as
\texttt{mech\_interp\_impossibility} in
\texttt{MechInterpInstance.lean} (three type axioms plus the
\texttt{circuitSystem} axiom encoding the Rashomon witness).

\paragraph{Critical framing: what this does and does not say.}
This result does \emph{not} undermine mechanistic interpretability---it
clarifies what MI can and cannot achieve.  Circuit discovery methods such as
activation patching~\citep{conmy2023automated}, sparse probing~\citep{elhage2022toy},
and circuit tracing~\citep{anthropic2025circuit} remain powerful tools for
understanding neural network internals.  What the impossibility establishes
is that the \emph{search for THE circuit}---a single, definitive
decomposition---is provably constrained: any method that commits to a
unique answer must sacrifice either faithfulness or stability.

The constructive path forward is to \emph{characterize the space of valid
circuits}, not to search for a single canonical one.  This is precisely
the direction that leading MI work is already moving toward:
\citet{anthropic2025circuit} report confidence levels rather than single
circuits, and \citet{conmy2023automated} frame automated circuit discovery
as search over a structured space.

\paragraph{Resolution: circuit equivalence classes.}
The uniform resolution of Section~\ref{sec:resolution} specializes to
mechanistic interpretability as follows.  Let $G = [C]$ denote the
equivalence class of circuits that reproduce the same I/O behavior.
The $G$-invariant explanation aggregates over all valid circuits:
\[
  \mathsf{exp}^*(\theta)
  \;=\;
  \{C \in \sH : C \text{ is a valid circuit for } f_\theta\},
\]
the set of all circuit decompositions consistent with the network's
behavior.  This map is \textbf{faithful in expectation} (it captures
all structure shared across valid circuits), \textbf{stable} by
construction (functionally equivalent networks have the same set of
valid circuits), but \textbf{not decisive}: it returns a set rather
than a single circuit.  The unresolvable queries---which specific
components ``really'' implement a computation when multiple valid
assignments exist---correspond precisely to the undirected edges in
the causal discovery analogy (Section~\ref{sec:instance-causal}).

Reporting the full equivalence class is the principled response to the
impossibility, trading decisiveness for faithfulness and stability, exactly
as predicted by Theorem~\ref{thm:universal}.  This reframes the MI
research agenda: the goal is not to find the one true circuit, but to
characterize which circuit features are \emph{invariant} across all valid
decompositions (the ``skeleton'' of the circuit space) and which are
decomposition-dependent (the ``reversible edges'').
